% !TeX program = xelatex
% ============================================================================
%  第五部分  多项式与生成函数
%
%  本部分由使用者另行提供，此处仅占位。保留以下内容的目的：
%    1. 让全书结构完整、可独立编译（XeLaTeX 两遍即可得到含第五部分的目录）；
%    2. 为后续补写标出章节编号、记号与交叉引用的落点。
%
%  补写时请注意：
%    - 章号从“第 1 章”开始（\\BeginPart 已把 chapter 计数器归零）；
%    - 超链接目标使用 \\theHchapter = <partno>.<chapter>，与前面四部分不会冲突；
%    - 记号沿用全书速查表：\\Z \\N \\Pp \\eps \\one \\id \\ord \\lcm \\rad \\iva
%     \\floor \\ceil，概率部分另有 \\E \\Var；若需新的宏，请加在 main.tex
%     的“简写命令”一节，不要在各部分文件里重复定义。
%    - 前置知识交叉引用：第三部分（组合恒等式、递推、格路）、第四部分
%     （概率生成函数、随机游走）、第二部分（模意义下的多项式与 NTT）。
% ============================================================================

\chapter{多项式与生成函数（占位）}
\begin{chapterintro}[章节导读]
本部分计划覆盖三块内容：\textbf{多项式运算与快速傅里叶变换}、\textbf{形式幂级数与普通生成函数}、\textbf{指数生成函数与指数公式}。它们是第三部分组合恒等式与递推的“机械化”版本，也是第一部分复数单位根的自然应用场所。此处仅保留骨架，正文由使用者提供。
\end{chapterintro}

\section{计划目录}
\begin{enumerate}
\item 多项式的表示与基本运算（系数形式、点值形式、乘法逆元）
\item 快速傅里叶变换与数论变换（NTT）
\item 形式幂级数：加法、乘法、逆、开方、对数与指数
\item 普通生成函数与线性递推的闭式（Binet 公式、有理生成函数）
\item 指数生成函数与指数公式、Bell 数与第二类 Stirling 数
\item 多项式与组合计数的综合应用
\end{enumerate}

\section{与全书的衔接点}
\begin{itemize}
\item \textbf{承接第一部分第 6 章}：单位根与复数乘法在那里引入，本部分用它构造 DFT。
\item \textbf{承接第三部分第 2、4 章}：二项式定理、范德蒙与朱世杰恒等式都是生成函数系数的特例；斐波那契、卡特兰数列的闭式由有理／代数生成函数给出。
\item \textbf{承接第三部分第 6 章}：格路计数的反射原理与生成函数是同一问题的两种语言。
\item \textbf{承接第四部分第 9、10 章}：概率生成函数 $G_X(z)=\E[z^X]$ 把离散分布的卷积变成多项式乘法；生日悖论的碰撞概率也可由生成函数重新推导。
\item \textbf{承接第二部分}：模素数意义下的 NTT、以及“多项式同余”与“整数同余”的类比。
\end{itemize}

\begin{tip}[给补写者的一句话]
生成函数的全部力量来自一件事：\textbf{把“数数”翻译成“多项式的系数”}。凡是第三部分里用分类讨论得到的递推，都可以在这里用一次乘法或一次求逆重新得到。补写时建议先复用第三部分的例题，再推广到新场景。
\end{tip}
